\begin{figure}[t]
\centering
\resizebox{\linewidth}{!}{%
\begin{tikzpicture}[
  font=\footnotesize, >=Stealth,
  inp/.style={draw=black!55, rounded corners=2pt, fill=black!4, align=center, text width=1.9cm, minimum height=0.62cm},
  base/.style={draw=black!45, rounded corners=2pt, fill=black!10, align=center, text width=1.95cm, minimum height=0.62cm},
  oursA/.style={draw=orange!80!black, very thick, rounded corners=3pt, fill=orange!14, align=center, text width=2.7cm},
  oursB/.style={draw=blue!70!black, very thick, rounded corners=3pt, fill=blue!10, align=center, text width=2.7cm},
  oursBs/.style={draw=blue!70!black, very thick, rounded corners=2pt, fill=blue!10, align=center, text width=1.95cm, minimum height=0.62cm},
  oursC/.style={draw=green!50!black, very thick, dashed, rounded corners=3pt, fill=green!8, align=center, text width=1.9cm},
  outn/.style={draw=black!55, rounded corners=2pt, fill=white, align=center, text width=1.8cm, minimum height=0.62cm},
  lab/.style={font=\scriptsize, fill=white, inner sep=1pt}]

\node[inp] (lang) at (1.2,6.7) {语言指令 $\ell$};
\node[inp] (img)  at (1.2,5.85) {图像 $I^{f}_t,\,I^{w}_t$};
\node[inp] (st)   at (1.2,5.0) {本体状态 $s_t$};
\node[oursC, minimum height=1.4cm] (c3) at (1.2,3.1) {\textbf{\Ciii{} 事件鲁棒训练}\\\textbf{与免标注}\\\textbf{真机适配}};
\node[inp, minimum height=0.9cm] (buf) at (1.2,1.3) {历史帧缓存\\$I_{0:t}$};

\node[oursA, minimum height=1.3cm] (pedet) at (5.0,3.9) {\textbf{\Ci{} 采摘事件检测器}\\\textbf{PEDet}\\{\scriptsize 腕部视觉 + 本体感知}\\{\scriptsize 因果 TCN，区分成败}};
\node[oursB, minimum height=1.8cm] (hem) at (5.0,1.3) {\textbf{\Cii{} 事件中心分层记忆}\\\textbf{HEM}\\{\scriptsize 场景锚 / 工作事件 /}\\{\scriptsize 情景摘要 / 时间上下文}};

\begin{scope}[on background layer]
\node[draw=black!40, dashed, rounded corners=4pt, fill=black!2, minimum width=8.6cm, minimum height=7.0cm] (vla) at (11.85,3.75) {};
\end{scope}
\node[font=\scriptsize, text=black!60] at (11.85,0.5) {基础 VLA 结构（灰色模块沿用，不作为本文贡献）};
\node[base] (tok)  at (8.9,6.7) {分词器};
\node[base, minimum height=0.75cm] (enc1) at (8.9,5.85) {视觉编码器\\{\scriptsize SigLIP（冻结）}};
\node[base, minimum height=1.6cm] (vlm) at (12.0,6.275) {VLM 骨干\\{\scriptsize Gemma-2B}};
\node[base, minimum height=1.2cm] (act) at (14.85,3.7) {动作专家\\{\scriptsize 流匹配}};
\node[base, minimum height=0.9cm] (mexp) at (14.85,1.3) {记忆分支};
\node[base, minimum height=0.9cm] (enc2) at (8.9,1.3) {视觉编码器\\{\scriptsize 共享，16 token/帧}};
\node[oursBs, minimum height=0.9cm] (emb) at (12.0,1.3) {\textbf{\Cii{} 类型-时间}\\\textbf{嵌入}};

\node[outn, minimum height=0.9cm] (chunk) at (17.6,3.7) {动作块 $A_t$\\{\scriptsize $20\times8$，执行 16 步}};
\node[outn] (guard) at (17.6,2.5) {安全包络};
\node[outn] (arm)   at (17.6,1.3) {机械臂 20~Hz};

\draw[->] (lang) -- (tok);
\draw[->] (img) -- (enc1);
\draw[->] (tok) -- (vlm.west |- tok);
\draw[->] (enc1) -- (vlm.west |- enc1);
\draw[->] (5.8,5.85) -- node[lab, right]{腕部} (5.8,4.55);
\draw[->] (st.east) -- (4.2,5.0) -- (4.2,4.55);
\draw[->, very thick, orange!80!black] (pedet) -- node[lab, right]{事件表 $\mathcal E_t$} (hem);
\draw[->] (buf) -- (hem.west |- buf);
\draw[->, very thick, blue!70!black] (hem.east |- enc2) -- node[lab, above]{$B{=}16$ 幅} (enc2);
\draw[->] (enc2) -- (emb);
\draw[->, very thick, blue!70!black] (emb) -- (mexp);
\draw[->] (vlm.east) -| (act.north);
\draw[->] (mexp) -- (act);
\draw[->] (act) -- (chunk);
\draw[->] (chunk) -- (guard);
\draw[->] (guard) -- (arm);
\draw[->, dashed, green!45!black, thick] (c3.east |- 0,3.6) -- node[lab, above]{伪标签} (pedet.west |- 0,3.6);
\draw[->, dashed, green!45!black, thick] ([yshift=0.25cm]c3.south east) -- (3.05,2.65) |- node[lab, pos=0.75, above]{扰动事件} (hem.west |- 0,1.9);

\node[font=\scriptsize, anchor=north] at (9.4,-0.05) {\sq{orange!80!black}\ \Ci{} 何时发生了采摘\qquad \sq{blue!70!black}\ \Cii{} 记住什么、如何编码\qquad \sq{green!45!black}\ \Ciii{} 如何对检测误差鲁棒、如何零标注上真机};
\end{tikzpicture}}
\caption{HarvestMem 总体框架。灰色为通用 VLA 的沿用结构：冻结的 SigLIP 编码当前图像，VLM 骨干融合指令（本体状态以离散化 token 一并输入，图中省略），流匹配动作专家输出 20 步动作块。彩色为本文创新：\Ci{} \pedet{} 从腕部图像和本体状态因果地检测成功采摘；\Cii{} \hem{} 依据事件表把 16 幅记忆的预算分配给场景锚、工作事件、情景摘要拼图帧和时间上下文，并在共享视觉编码器之后叠加类型-时间嵌入；\Ciii{} 训练时按检测器误差模型扰动事件，部署时由作业结果事后确认生成伪标签、持续更新 \pedet。记忆图像数和 token 数与基础结构完全相同。}
\label{fig:overview}
\end{figure}
